% Generated draft: requires mathematical and source review.
\documentclass[11pt]{article}
\usepackage[margin=25mm]{geometry}
\usepackage{fontspec}
\setmainfont{Latin Modern Roman}
\usepackage{amsmath,amssymb,mathtools,unicode-math}
\setmathfont{Latin Modern Math}
\usepackage{enumitem,needspace,xurl,hyperref}
\setlength{\parindent}{0pt}
\setlength{\parskip}{5pt}
\setlength{\emergencystretch}{3em}
\newcommand{\N}{\mathbb{N}}
\newcommand{\Z}{\mathbb{Z}}
\newcommand{\Q}{\mathbb{Q}}
\newcommand{\R}{\mathbb{R}}
\newcommand{\C}{\mathbb{C}}
\newcommand{\F}{\mathbb{F}}
\newcommand{\A}{\mathbb{A}}
\newcommand{\PP}{\mathbb{P}}
\newcommand{\E}{\mathbb{E}}
\newcommand{\eps}{\varepsilon}
\newcommand{\dd}{\,\mathrm{d}}
\begin{document}
\section*{UnsolvedMath ID: 30003678}
\textbf{Source problem code:} catalogue entry
\subsection{Definitions and mathematical statement}
Let \(d\in\N\) with \(d\geq3\). The supplied statement concerns a scale-free percolation graph \(G=(\Z^d,E)\) in a regime having a unique infinite connected component \(\mathcal C_\infty\), with parameters satisfying
\[
\gamma>2,\qquad \alpha>2d.
\]
The inequalities are strict; their boundary cases are not included. The record does not specify the graph's probability law, the meanings of \(\alpha\) and \(\gamma\), or their relation to any weight distribution or connection probabilities. These inequalities therefore describe the supplied regime but do not define the model.

Provided that \(\mathcal C_\infty\) is locally finite, write \(\deg_{\mathcal C_\infty}(x)\) for the number of neighbors of \(x\) in this component. For a fixed realization of the graph, simple random walk \((X_n)_{n\geq0}\) on \(\mathcal C_\infty\) has transition probabilities
\[
\mathbb P^G(X_{n+1}=y\mid X_n=x)
=
\begin{cases}
\deg_{\mathcal C_\infty}(x)^{-1},&\{x,y\}\in E,\\
0,&\text{otherwise}.
\end{cases}
\]
Here \(x,y\in\mathcal C_\infty\), and \(\mathbb P_x^G\) denotes the walk law started at \(x\). Define
\[
T_x^+=\inf\{n\geq1:X_n=x\},\qquad \inf\varnothing=\infty.
\]
The walk is recurrent if \(\mathbb P_x^G(T_x^+<\infty)=1\) for every \(x\in\mathcal C_\infty\), and transient if this probability is strictly less than \(1\) for every such \(x\). On a connected, locally finite graph these alternatives do not depend on the starting vertex. Local finiteness is needed for this transition rule and is not established by the supplied record.

\textbf{Requested mathematical target.} Determine whether simple random walk on the unique infinite cluster is recurrent or transient in the stated scale-free percolation regime, for dimensions \(d\geq3\) and the admissible parameter choices. The record does not explicitly specify the probability quantifier over graph realizations; an almost-sure classification would require the missing model law and a specification of whether the infinite-cluster assumption holds almost surely or is imposed by conditioning.

The source also asks for the analogous classification for long-range percolation. Its original wording refers to ``regime (ii) of the stated theorem,'' but neither that theorem nor the long-range percolation law and parameter regime is supplied. Consequently, a precise long-range percolation assertion, and a complete model-specific scale-free percolation assertion, cannot be recovered from this record alone.
\subsection{Short English statement}
For scale-free percolation on \(\Z^d\), with \(d\geq3\), a unique infinite cluster, \(\gamma>2\), and \(\alpha>2d\), determine whether simple random walk is recurrent or transient. The source also asks the analogous long-range percolation question. Missing model definitions and the omitted theorem prevent a fully specified formulation.
\subsection{Sources}
\begin{itemize}
\item Local registry: \texttt{data/unsolvedmath/problems.json}, permanent UnsolvedMath ID \(30003678\), source problem code \texttt{catalogue entry}, title ``Recurrence of High-Dimensional Scale-Free Percolation.'' Attribution is to an Oberwolfach workshop open-problem discussion; no individual proposer is supplied.
\item Supplied source citation: \emph{Network Models: Structure and Function (2018)}, Oberwolfach Reports, pp.~3427--3470, DOI \texttt{10.4171/owr/2017/57}. \url{https://doi.org/10.4171/owr/2017/57}. The registry records proposed year 2017. These bibliographic dates are retained as supplied.
\item Supplied literature reference: \emph{Scale-free percolation and long-range percolation}. \url{https://doi.org/10.1214/16-AAP1270}.
\item Supplied literature reference: \emph{Continuity and shape theorem for long-range percolation}. \url{https://arxiv.org/abs/2312.04099}.
\item Supplied literature reference: \emph{OWR Report 57/2019}. \url{https://ems.press/content/serial-article-files/46721}.
\item Literature assessment dated exactly \texttt{2026-08-21}: the supplied curation labels the problem open, describes the stated strong-decay scale-free percolation question as unresolved, and says that later long-range percolation advances do not settle that exact model question. This is supplied dated curation, not independently verified here; its time-sensitive status and cited claims require independent verification.
\item The registry labels the statement \texttt{corrected\_verified} and reports that the official Oberwolfach report was checked directly and that the cleaned statement restored definitions while preserving the target. These are supplied verification claims, not verification performed for this draft; the complete model definitions are absent from the supplied record.
\end{itemize}
\end{document}
